\section{Functional linkage and replication}
\label{sec:linkage-replication}

We now show that the dimension of the active response space is additive
under generic direct sums.  The main point is to turn the linear relations
between linked responses into identities between eigenvalue branches.
These identities will continue to control responses through partners of
every allowed size.

Retain the progression $\mathcal S=\{pN:N\geq 1\}$ and the common
eigenvalue multiplicity $t$.  Thus, at size $n\in\mathcal S$, the number
of distinct generic eigenvalues is $e_n=dn/t$.  We use the following
previously established facts.  At a direct sum of independent small
generic inputs of sizes in $\mathcal S$, the spectra are disjoint and
have the maximal number of distinct eigenvalues.  Their commuting
Taylor deformations therefore have scalar spectral blocks, by
Lemma~\ref{lem:scalar-branches}.  The response identities admit the
specializations of Lemma~\ref{lem:active-specialization}, and their
leading terms satisfy Proposition~\ref{prop:word-transfer}.

For a generic input $P$, let $\Omega_P$ be its spectral splitting
field, and choose a constant input direction $v$ such that
$D_h(P;v)\neq 0$ for every eigenvalue label $h$.  Recall the active-space
partition $\mathcal P_P$, characterized by
\begin{equation}
 \mathcal T_P
 =\left\{c\in\Omega_P^{e_n}:
       \frac{c_h}{D_h(P;v)}
       \text{ is constant on each member of }\mathcal P_P\right\}.
 \label{eq:linkage-partition}
\end{equation}
Labels in one member of this partition are called linked.  In particular,
$a_n=\dim_{\Omega_P}\mathcal T_P$ is the number of its members.
The argument below uses the full coefficientwise definition of
$\mathcal T_P$: it contains the loop responses $D(P;Y)$ and the
coefficient vectors of $m(P,Q;U,V)$ for every partner size in
$\mathcal S$.

\subsection{A formal identity for linked branches}

\begin{proposition}[Functional linkage]
\label{prop:functional-linkage}
Let $P$ be a small generic input of size $p$, and let $h,h'$ be linked
labels at $P$.  Let $C_1,\ldots,C_s$ be small generic inputs of sizes in
$\mathcal S$, jointly independent with $P$, where $s=0$ is allowed.
Put $R=\operatorname{diag}(P,C_1,\ldots,C_s)$.
Over a common spectral splitting field $E$, denote by
$\Lambda_j^R(\Delta)$ the scalar eigenvalue branch of $W(R+\Delta)$
whose value at zero is $\lambda_j(P)$, for $j=h,h'$.

Define
\[
 g_j(\eta)=\lambda_j(P+\eta v)-\lambda_j(P)
 \in\Omega_P[[\eta]],\qquad j=h,h'.
\]
Then $g_{h'}'(0)\neq 0$, and
\begin{equation}
 \Lambda_h^R(\Delta)-\lambda_h(P)
 =\bigl(g_h\circ g_{h'}^{-1}\bigr)
       \bigl(\Lambda_{h'}^R(\Delta)-\lambda_{h'}(P)\bigr).
 \label{eq:functional-linkage}
\end{equation}
Here $g_{h'}^{-1}$ denotes the inverse for composition.  In particular,
the univariate series in this identity depends only on $P,h,h'$ and $v$,
and is unchanged when auxiliary blocks are adjoined.
\end{proposition}

\begin{proof}
We first explain why the response relations persist while $P$ moves.
By the Taylor convention in Section~\ref{subsec:formal-conventions},
substitution $P\mapsto P+\eta v$ gives homomorphisms
\[
 \operatorname{Frac}K[[P]]
 \longrightarrow \bigl(\operatorname{Frac}K[[P]]\bigr)[[\eta]],
 \qquad
 \Omega_P\longrightarrow\Omega_P[[\eta]],
\]
both reducing to the identity at $\eta=0$. The second uniquely
extends the first through the finite separable extension
$\Omega_P/\operatorname{Frac}K[[P]]$. Thus every nonzero element
has image with invertible constant term, and the eigenvalues and
projections map to their uniquely labeled Taylor deformations.

For linked labels, Equation~\eqref{eq:linkage-partition} gives
\begin{equation}
 D_{h'}(P;v)c_h-D_h(P;v)c_{h'}=0
 \label{eq:linked-response-relation}
\end{equation}
for every loop response and every coefficient vector of every paired
response.  Apply the preceding homomorphism coefficientwise in the
partner coordinates.  The resulting identities are exactly
Equation~\eqref{eq:linked-response-relation} at the moved input
$P+\eta v$, with the moved eigenvalues, projections, and $D$-factors.
They remain valid in the triangular specializations used in
Proposition~\ref{prop:word-transfer}.  More explicitly, first clear the
finitely many auxiliary resolvent determinants.  At each Taylor degree,
the number of strictly upper steps is bounded.  This permits their
entries to have Laurent coefficients in $\eta$, while the diagonal
entries are expanded at the original small inputs.  Coefficientwise
convergence and specialization are then supplied by
Lemma~\ref{lem:active-specialization}.  Thus the transfer argument
applies over $E((\eta))$ with the transported response relation.

Write $v_R$ for $v$ supported in the $P$-block and put
$R(\eta)=R+\eta v_R$.  At $R(\eta)+H$, impose the equation
\begin{equation}
 \Lambda_{h'}^{R(\eta)}(H)=\lambda_{h'}(P+\eta v).
 \label{eq:fixed-linked-branch}
\end{equation}
Choose linear coordinates $H=u v_R+H_{\mathrm{rest}}$, with
$H_{\mathrm{rest}}$ in a fixed complementary space of input directions.
The derivative of the left side in $u$ at $H=0$ is
$D_{h'}(P+\eta v;v)$, a unit in $E[[\eta]]$.  Formal implicit
substitution therefore eliminates $u$.  In particular, the quotient
of $E[[\eta,H]]$ by Equation~\eqref{eq:fixed-linked-branch} is
\begin{equation}
 E[[\eta,H_{\mathrm{rest}}]],
 \qquad
 E[[\eta,H_{\mathrm{rest}}]]
 \lhook\joinrel\longrightarrow
 E((\eta))[[H_{\mathrm{rest}}]].
 \label{eq:linkage-quotient-injection}
\end{equation}
The eliminated series has zero constant term in $H_{\mathrm{rest}}$.

On this quotient, scalar-block separation implies that the fixed-root
Schur expression $S_{h'}$ at $R(\eta)$ vanishes.  We claim that $S_h$
also vanishes over the right-hand ring in
Equation~\eqref{eq:linkage-quotient-injection}.  Induct on total degree
in $H_{\mathrm{rest}}$.  Both Schur expressions vanish in degree zero.
If they vanish below degree $l$, the commuting-parameter conclusion of
Proposition~\ref{prop:word-transfer}, together with the transported
relation~\eqref{eq:linked-response-relation}, gives scalar leading
coefficients $[S_j]_l=c_{j,l}\pi_j$ for $j=h,h'$, with
\[
 D_{h'}(P+\eta v;v)c_{h,l}
   =D_h(P+\eta v;v)c_{h',l}=0.
\]
Here the $c_{j,l}$ are scalar homogeneous polynomials in the remaining
parameters over $E((\eta))$; the projections are the moved target
projections. The first factor is nonzero, so $c_{h,l}=0$ and
$[S_h]_l=0$.

It follows that the $h$-branch is fixed as well:
\[
 \Lambda_h^{R(\eta)}(H)=\lambda_h(P+\eta v)
\]
on the quotient.  For clarity, the implication from $S_h=0$ uses the
scalar blocks of Lemma~\ref{lem:scalar-branches}: the complementary
block at the fixed root is invertible, and the determinant of the full
matrix at that root is a unit times the $t$-th power of the shift of
the $h$-branch.  The quotient is a domain, so the shift is zero.
Both branch series and the eliminated coordinate are regular in
$\eta$.  The injection in
Equation~\eqref{eq:linkage-quotient-injection} therefore descends this
equality to $E[[\eta,H_{\mathrm{rest}}]]$.

Finally make the invertible linear coordinate change
$\Delta=H+\eta v_R$.  Equation~\eqref{eq:fixed-linked-branch} becomes
\[
 \Lambda_{h'}^R(\Delta)-\lambda_{h'}(P)=g_{h'}(\eta).
\]
Since $g_{h'}'(0)=D_{h'}(P;v)\neq0$, it eliminates $\eta$ as
$g_{h'}^{-1}(\Lambda_{h'}^R(\Delta)-\lambda_{h'}(P))$.
Substitution in the equality for the $h$-branch gives
Equation~\eqref{eq:functional-linkage}.
\end{proof}

\subsection{All partner sizes and the dimension formula}

We next pass from an identity near a direct sum to an identity for
every coefficient defining the generic active space.  This step is
necessary because a bound proved using a fixed collection of partners
would not bound $a_n$.

\begin{proposition}[Replication of active dimensions]
\label{prop:active-replication}
For every positive integer $N$,
\[
 a_{pN}=N a_p.
\]
\end{proposition}

\begin{proof}
Put $n=pN$ and choose independent small generic inputs
$P_1,\ldots,P_N$ of size $p$.  Write
$R_1=\operatorname{diag}(P_1,\ldots,P_N)$ and use
$x=R_1+\Delta_1$ to label the generic size-$n$ spectrum by its branches
at $R_1$.

This substitution embeds $K[[x]]$ in $E[[\Delta_1]]$.  Indeed, for
an image of a series over $K$, the Taylor coefficients are uniformly
integral in the small parameter $\varepsilon$.  Substituting
$\Delta_1=\varepsilon X-R_1$, with independent coordinates $X$,
is therefore convergent and recovers the original series at
$\varepsilon X$.  Homogeneous expansion at that point detects every
nonzero series.  The embedding extends to the fraction field.  All
generic spectral roots are supplied by the scalar branches at $R_1$,
so it also embeds a spectral splitting field $\Omega_x$ into
$\operatorname{Frac}E[[\Delta_1]]$, with these labels.  The branches
and spectral projections themselves are regular at $\Delta_1=0$.

Fix a block $i$ and labels $h,h'$ in one member of
$\mathcal P_{P_i}$.  Choose a constant direction $v_i$ supported in
that block, nonzero under every $D_h(P_i;\cdot)$.  Apply
Proposition~\ref{prop:functional-linkage} first at $R_1$.
Differentiating its identity gives proportional loop responses at
$x$.  Now adjoin an independent generic partner $Q$ of any fixed
size in $\mathcal S$, and apply the same proposition to
$\operatorname{diag}(R_1,Q)$.  Its univariate series is the same
one.  At $\operatorname{diag}(x,Q)$, perturb the opposite input
blocks by $uU$ and $zV$.  Either arrow alone preserves the spectrum,
so the first derivatives in $u$ and $z$ vanish.  The coefficient of
$uz$ in the branch is the paired response $m_h(x,Q;U,V)$.
The mixed chain rule and the separate first derivative in $v_i$
give
\begin{equation}
 D_{h'}(x;v_i)m_h(x,Q;U,V)
   -D_h(x;v_i)m_{h'}(x,Q;U,V)=0.
 \label{eq:replicated-pair-relation}
\end{equation}
The two $D$-factors have nonzero reductions at $\Delta_1=0$.

We must verify this relation coefficientwise in $Q$, rather than only
at its generic value.  Cross multiplication has already removed the
response ratio.  Multiply further by the finitely many determinants
of $W(Q)-\lambda_h(x)$ and $W(Q)-\lambda_{h'}(x)$ needed to clear
the auxiliary inverses.  At each fixed $\Delta_1$-degree, the resulting
series in the independent coordinates of $Q$ has coefficients
uniformly bounded below in $\varepsilon$-valuation.  The branch and
projection coefficients do not involve $Q$; the other factors are
ordinary matrix series with two prescribed arrow steps and a fixed
number of Taylor insertions.  Their path expansions give precisely
this bound, as in Lemma~\ref{lem:admissible-evaluation}.
Evaluation at small generic $Q$ therefore detects every coefficient.
Equation~\eqref{eq:replicated-pair-relation} is consequently a
formal identity in those coordinates.  Injectivity of the embedding
of $\Omega_x$ descends it to the generic size-$n$ field.  Finally,
the cleared determinants have nonzero constant terms in $Q$, so
they can be canceled in $\Omega_x[[Q]]$.

This argument applies separately to every partner size with the same
labeling of $x$ and the same $D$-factors.  It therefore gives a
relation on every coefficient vector used to define $\mathcal T_x$.
Together with the corresponding loop relations, it says that
$c_h/D_h(x;v_i)$ is constant on each old partition member within
block $i$, for every $c\in\mathcal T_x$.  These partitions have
$Na_p$ members in total, and all their normalizing factors are
nonzero.  Thus $a_n\leq Na_p$.

For the opposite inequality, use
Lemma~\ref{lem:active-specialization} for each $P_i$, choosing all
partners jointly independent.  At $R_1$, support each loop direction,
and both arrows of each paired response, in its chosen block $i$.
The responses then coincide with the original responses on that
block and vanish on every other block.  They have rank $Na_p$.
Evaluate the same finite collection first at $R_1+\Delta_1$.
All branch data and partner resolvents are regular at
$\Delta_1=0$, since the partner spectra are disjoint from that of
$R_1$.  A generic rank bound smaller than $Na_p$ would make every
minor of that order vanish over the embedded generic field, and hence
also at the reduction.  The displayed specialization contradicts this.
Therefore $a_n\geq Na_p$.
\end{proof}

\subsection{Partners of one common size}

The dimension formula permits us to use a fixed matrix size on both
ends of every paired response.  We will also need to fix one arrow
and vary only its reverse.

\begin{corollary}[Equal-size partners]
\label{cor:equal-size-partners}
There are $r\in\mathcal S$ and a positive integer $K_0$ with the
following property.  Let $P,Q_1,\ldots,Q_{K_0}$ be independent small
generic inputs of size $r$.  For generic fixed input arrows
$U_j:P\leftarrow Q_j$, the loop columns and the images of
\[
 V\longmapsto m(P,Q_j;U_j,V),
 \qquad V:Q_j\leftarrow P,\quad 1\leq j\leq K_0,
\]
span $\mathcal T_P$ after extension to the joint field of the data.
The analogous assertion holds with the other arrow fixed.  The same
$r$ and $K_0$ may be used for both assertions.
\end{corollary}

\begin{proof}
Choose a finite witness list at size $p$ from
Lemma~\ref{lem:active-specialization}, and let its partner sizes be
$q_1,\ldots,q_M\in\mathcal S$.  Choose $r=pN$ with
$r\geq q_j$ for every $j$.  Specialize the size-$r$ base to a direct
sum of $N$ independent size-$p$ inputs.  For each block, use a fresh
copy of the witness list.  Pad a size-$q_j$ partner to size $r$ by
adjoining independent generic diagonal data of size $r-q_j$, omitting
the padding if this difference is zero.  Give both arrows zero entries
on the padding.  The response remains the original size-$p$ witness
on its chosen base block and is zero on the other base blocks.

The resulting columns have rank $Na_p=a_r$ by
Proposition~\ref{prop:active-replication}.  All inputs at this
specialization have maximal spectral count, and the partner spectra
are disjoint from the base spectrum.  The regular-minor argument in
the preceding proof therefore transfers this rank to generic
size-$r$ base and partner inputs.

Assign a separate partner to every paired column in this finite
list.  At the witness specialization, choose its forward arrow to
be the specified $U_j$ and test the variable reverse arrow on the
specified $V_j$.  A minor of order $a_r$ is nonzero there.  Since
paired responses are bilinear in their arrows, this minor remains
nonzero for generic fixed forward arrows and suitable reverse-arrow
columns.  This proves the first assertion with finitely many partners.
The identical argument fixes the reverse arrows instead.  Increasing
the number of independent partners to the larger of the two finite
numbers proves both assertions with one $K_0$.  If no paired column
is needed, one unused generic partner may be added so that $K_0>0$.
\end{proof}
